\section{Systems Realization and the Four-Arm Measurement}
\label{sec:systems}

The architectural claim is meaningful only if the runtime enforces the
no-escape/no-store inference contract (Figure~\ref{fig:lifecycle},
Appendix~\ref{app:lifecycle}). Production realization requires generator
placement adjacent to the FFN kernel emitting core values into kernel-local
scratch; associative execution ($XU_r$, core apply, $ZV_r$); a cheap
eligibility router in the runtime path; and lifecycle enforcement that
returns only the FFN output while instrumenting the core address range to
verify no global store, alias, state registration, or unselected-branch
launch. If hardware lacks software-managed SRAM, a cache-resident
implementation may still be useful but cannot claim the strict contract
without an address trace showing the generated region never reaches HBM.

\paragraph{Four-arm protocol.}
Isolating lifetime from low rank requires the four arms of
Table~\ref{tab:four-arms} (Appendix~\ref{app:fourarms}): arm 3 and arm 4
use the same generator, rank, core values, precision, and output tolerance,
differing only in whether the generated core is materialized to global
memory (arm 3) or consumed fused from scratch without a store (arm 4). The
historical benchmark arm (Appendix~\ref{app:histtables}) is a persistent
low-rank forward--backward microbenchmark whose speedup does \emph{not}
isolate ephemerality---which is precisely why the fused arm-4 measurement
was required.


\paragraph{Measured result.}
Arm 4 has been \emph{executed and measured on device}
(Table~\ref{tab:fourarmmeasured}): an audited Ascend-C operator (a static
source audit mechanically enforces the no-store contract; a test suite checks
the five \texttt{GlobalTensor}s, the single output write, the absence of any
core \texttt{DataCopy}, and deterministic execution) was validated for
algebraic equivalence against arm 3 under a hash-bound contract manifest
(\texttt{max\_abs} $0.00195$, at the fp16 rounding floor) and instrumented
with matched \texttt{msprof} counter passes at two geometries. The
\texttt{ephemeral} kernel is visible in the device trace; the arm-3-only core
round-trip ops (\texttt{Tanh}/\texttt{Muls}/\texttt{Adds}/\texttt{Cast} to
global memory) are absent in the fused trace---direct counter evidence that
the core does not escape. The first configuration lost to materialization
($1.228$ vs.\ $0.847$~ms) entirely because per-call host work (tiling
generation, workspace allocation, fp32 casts) dominated the loop, not device
execution. After host fast-pathing (cached tiling/workspace, cached fp16
weight image), the fused arm reaches \emph{parity} with materialization at
the probe geometry. Across repeated independent same-session runs the fused
and materialized medians stay within a $\pm3\%$ band straddling $1.0\times$
(the final shipped confirmation measures fused $0.823$~ms vs.\ materialized
$0.851$~ms, i.e.\ $3.2\%$ faster; Table~\ref{tab:fourarmmeasured} row
4$^\dagger$): once host work is amortized, there is no residual device-side
penalty for generating and consuming the ephemeral core on chip at
L2-resident geometries.

The advantage inverts at a $T{=}4096$ geometry.  Note the two memory
hierarchies separately: the logical per-projection controller output
$C\in\R^{T\times r}$ is $2$~MiB in fp32 at $T{=}4096$, $r{=}128$ ($6$~MiB
across the three projections if all were live simultaneously; v1 processes
them one at a time), and it exceeds the $192$~KiB
per-core Unified Buffer at every geometry (this is why v1 tiles), while
L2 residency is not claimed from L2 counters (which read zero under this
\texttt{msprof} build) but inferred from \texttt{Memory}-pass DRAM byte
totals (Appendix~\ref{app:counters}); the exact L2 capacity of this
910B2 is not separately measured here, and UB overflow is not evidence
of L2 overflow.  At $T{=}4096$ equivalence still
passes, but the fused arm is $1.93\times$ slower
($2.682$ vs.\ $1.393$~ms), with matched profiler passes localizing the
deficit to the v1 kernel's compute-bound single-core epilogue, not memory
traffic (Figure~\ref{fig:opdecomp}, Appendix~\ref{app:fourarmfig}).  This
deficit is a critical limitation of the current prototype, not a minor
tuning gap.

\paragraph{Architectural root cause.}
The v1 kernel runs the entire pipeline---cube matmul
($X_{\mathrm{ctrl}} W_{\mathrm{head}}^\top$, producing $C\in\R^{T\times r}$ in fp32),
elementwise vector epilogue, and output write---on a \emph{single} AICore
with \texttt{block\_dim}$=1$.  The kernel tiles $T$ at
\texttt{BASE\_M}$=128$ and holds one $128\times r$ fp32 $C$ tile
($64$~KiB at $r{=}128$) in the 192~KiB Unified Buffer (UB) at a
time---the full $C$ is never resident---but at this tile size a second
$C$ buffer would not fit alongside activation staging, output, and
matmul scratch, so v1 cannot double-buffer: within each tile the
cube$\to$vector dependency is strictly serial (\texttt{GetTensorC} must
complete before \texttt{Tanh} begins), and cube and vector work never
overlap.  The vector epilogue alone
consumes more wall time than all library matmuls of either arm combined,
and the seven idle AICores cannot assist: the no-store contract requires
the core to be produced and consumed in one kernel invocation, but the
single-core serial pipeline cannot exploit the available parallelism
(sizing detail in Appendix~\ref{app:v4design}).

\paragraph{v4 multi-core redesign (summary; full design in
Appendix~\ref{app:v4design}).}
An earlier v3 attempt tried N-distribution (splitting the rank dimension
$r$ across cores) but triggered aicore exception 507015 under sequential
launches because the KFC matmul workspace was shared across cores without
isolation; it is archived as v3.  The redesigned v4 kernel addresses both
the v1 bottleneck and the v3 failure mode through three coordinated
mechanisms: \emph{M-only multi-core distribution} (each AICore processes a
disjoint token range, no cross-core synchronization), \emph{per-core
workspace isolation} (a unique, non-overlapping
\texttt{REGIST\_MATMUL\_OBJ} workspace slice per core, eliminating the v3
collision), and a \emph{double-buffered UB pipeline}
(\texttt{BUFFER\_NUM=2}, overlapping cube and vector work within each core
at a ${\sim}160$~KiB per-tile footprint within the 192~KiB budget).  The
design projected ${\sim}5$--$7\times$ over v1 at $T{=}4096$; the measured
decisive run (2026-10-08, single Ascend~910B2, hardened two-phase-ref
harness, 30 warmup/200 iterations/5 repeats per config) confirms the
direction but not the magnitude: v4 is algebraically identical to v1
($\max|\Delta_{v1-v4}|=0$ at every config) and $1.13\times$, $1.38\times$,
$1.71\times$, $2.14\times$, and $3.03\times$ faster at $T{=}128$, $512$,
$1024$, $2048$, and $4096$ respectively---clearing the ${\sim}2\times$
viability threshold at $T\ge2048$ but falling short of the projection,
consistent with partial L2-contention saturation from shared-$B$ reads
(remaining risks: UB capacity would overflow at $r\ge256$ without further
tiling, and the 1000-call stability check at $T{=}4096$ hit an
environment-level NaN transient in its unhardened comparison path).
The no-store contract is preserved: each core computes the matmul and
epilogue entirely in UB and writes only the final fp16 output to global
memory.  The v4 kernel source and benchmark harness are archived in
\texttt{custom\_ops/ephemeral\_core\_apply/torch\_ext/v4\_multicore/}.

\begin{table}[t]
\centering
\caption{Measured four-arm ladder on a single Ascend 910B2 (CANN 8.5,
torch\_npu 2.9), geometry $(d,m)=(4096,14336)$, rank 128, fp16, seed 1234;
host timer with device synchronize. Arm 4 runs the audited, hash-bound
Ascend-C operator; equivalence against arm 3 passed at $\max|\Delta|=
1.95\times10^{-3}$ (fp16 rounding floor; tolerance $5\times10^{-3}$).
Row 4$^\dagger$ adds the host fast path (cached tiling/workspace, cached
fp16 weight image), leaving device execution identical; it is from a
separate confirmation run, so the fused-versus-materialized comparison
uses same-run pairs: the row-4$^\dagger$ run's matched materialized
median was $0.8509$~ms (same-run delta $-3.2\%$), and the shipped fused
arm sits \emph{at parity} with the materialized arm (medians within a
$\pm3\%$ band straddling $1.0\times$). Full timing protocol in
Section~\ref{sec:protocol}.}
\label{tab:fourarmmeasured}
\begin{tabular}{llrr}
\toprule
Arm & Execution & Median ms & $\pm$ std \\
\midrule
1 & Dense kernels & 0.4535 & 0.0021 \\
2 & Persistent low-rank & 0.2469 & 0.0016 \\
3 & Generated core, global-memory materialization & 0.8689 & 0.0312 \\
4 & Fused no-store (v1, per-call host setup) & 1.2315 & 0.0315 \\
4$^\dagger$ & Fused no-store (host fast path, shipped) & 0.8234 & 0.0271 \\
\bottomrule
\end{tabular}
\end{table}

Table~\ref{tab:positioning} (Appendix~\ref{app:positioning}) places this
measured status against the closest mechanism-level comparisons.
